%% 第一章

\chapter{绪论}
\chaptermark{绪论}

\section{研究背景及意义}

大语言模型（Large Language Models, LLMs）改变了人类获取知识与完成任务的方式。从2018年BERT\cite{devlin2018bert}和GPT\cite{radford2018improving}开始，大语言模型的参数规模从数亿增长到数千亿。近年来，以GPT-5.5\cite{openai2026gpt55}、Gemini 3.1 Pro\cite{geminiteam2026gemini31}、Claude Opus 4.7\cite{anthropic2026claudeopus47}为代表的国际前沿模型，以及DeepSeek-V4\cite{deepseekai2026deepseekv4}、Qwen3.6系列\cite{qwenteam2026qwen36}等国产大模型，在数学推理、代码生成、多语言理解等任务上表现出较强能力，应用范围仍在扩展。这些模型已经从学术研究走向商业部署，在教育、医疗、法律、科研等领域投入使用。

然而，模型能力提升和应用范围扩展并不意味着其输出天然可信。当前大语言模型的预训练仍主要依赖最大化下一词预测概率，该目标能够增强语言建模能力，却并不直接约束模型输出是否符合人类真实需求。因此，未经对齐的模型仍可能产生事实性幻觉（hallucination）\cite{ji2023survey}、有害内容及偏见性输出，在教育、医疗、法律等高风险场景下甚至可能造成更严重后果。仅依靠扩大模型规模和提升通用能力，难以保证模型在价值观、事实准确性与安全性上与人类期望一致。如何缓解预训练目标与人类真实需求之间的错位，使模型输出更加可靠、可控和安全，已成为大语言模型研究中的关键问题，即大模型对齐问题\cite{gabriel2020artificial}。

应对对齐问题的一类重要做法，是在后训练（post-training）阶段将人类期望转化为可优化的偏好监督或奖励信号，并利用该信号引导策略模型更新。典型的基于人类反馈的强化学习（Reinforcement Learning from Human Feedback, RLHF）方法通常先训练奖励模型，再用奖励模型评分指导策略优化。然而，奖励信号的可靠性会显著影响策略优化方向：偏好标注噪声可能使奖励模型学习到有偏代理目标，而固定奖励模型又难以适应策略优化过程中持续变化的生成分布。当生成分布逐渐偏离奖励模型训练分布时，奖励评分可能不再可靠，进而出现奖励模型过优化（reward model overoptimization）\cite{gao2023scaling}。因此，对齐训练不仅取决于策略优化算法，也取决于奖励信号能否在策略持续变化时保持可靠和稳定。

逆强化学习（Inverse Reinforcement Learning, IRL）\cite{ng2000algorithms}提供了另一种奖励学习视角：可以将高质量示范响应视为偏好目标在行为层面的体现，并从示范行为中反演可用于策略优化的奖励函数。相比依赖额外成对偏好比较训练奖励模型的方法，这一思路有助于减少对偏好比较监督的依赖。若奖励函数能够随策略生成分布进行迭代更新，该思路有助于缓解固定奖励模型带来的分布偏移；同时，逆强化学习也引入了新的问题：奖励函数一旦随策略反复更新，其数值尺度和语义含义都可能发生漂移，进而影响训练稳定性。

因此，本文将“大模型偏好对齐”这一较宽泛的问题进一步收束为一个具体的奖励学习问题：如何从高质量示范响应中反演出能够刻画示范响应所蕴含偏好目标的奖励函数，并在奖励函数随策略更新而变化时保持其尺度和语义稳定。本文中的高质量示范响应主要指人工或 AI 生成的高质量参考回答，其被视为偏好目标在行为层面的近似体现。

围绕上述问题，本文提出近端逆奖励优化方法（Proximal Inverse Reward Optimization, PIRO）。该方法在策略—奖励交替更新过程中引入近端约束，限制相邻迭代轮次中奖励函数的变化幅度，使奖励函数既能够随当前策略分布动态调整，又避免多轮训练中的尺度膨胀和语义漂移。由此，本文试图在奖励函数的适应性与稳定性之间取得平衡，为多轮逆强化学习式偏好对齐提供更加稳定的奖励学习机制。

\section{研究现状}

围绕大语言模型偏好对齐问题，相关研究可以从奖励信号的建模方式及其与策略更新的关系两个角度进行梳理。现有方法大体可分为三类：一类显式训练奖励模型，并将其作为后续策略优化的外部评分器；一类不单独训练奖励模型，而是直接将偏好关系转化为策略优化目标；还有一类将奖励函数作为可更新变量，在策略--奖励交替优化过程中动态学习奖励信号。不同路线在训练复杂度、奖励可解释性、分布适应性和工程实现难度上各有侧重，也不同程度地面临偏好数据噪声、固定奖励失配和动态奖励漂移等问题。

\subsection{基于人类反馈的强化学习}

基于人类反馈的强化学习（Reinforcement Learning from Human Feedback, RLHF）是大语言模型对齐中最具代表性的技术路线。其基本思路是利用人类偏好数据训练奖励模型，再以奖励信号引导策略优化。Christiano等人\cite{christiano2017deep}较早在强化学习任务中系统使用人类反馈，Ziegler等人\cite{ziegler2019fine}和Ouyang等人\cite{ouyang2022training}进一步将这一思路用于语言模型后训练，使RLHF成为当前大模型对齐的重要方法。

近几年，RLHF相关研究主要围绕训练效率、策略约束和在线数据更新展开：一方面，相关工作分析了PPO实现细节、奖励归一化和策略约束等因素对训练稳定性的影响\cite{zheng2023secrets,moskovitz2023fool}；另一方面，ReMax、GRPO等方法尝试降低价值函数依赖和采样开销，使强化学习式后训练更容易扩展到大模型场景\cite{li2023reward,shao2024deepseekmath}。此外，在线RLHF和基于当前策略样本的偏好学习强调在训练过程中引入当前策略样本和迭代反馈，以缓解固定离线数据难以覆盖持续变化模型分布的问题\cite{tajwar2024preference,online_rlhf_2025}。

上述改进提高了RLHF的工程可用性，但其局限仍主要来自奖励信号本身。一方面，成对偏好标注依赖标注者对不同回答进行比较，不同标注者的价值标准、任务理解和风险偏好可能并不一致，奖励模型因此可能学习到带有噪声和偏差的代理目标。另一方面，奖励模型通常在有限偏好数据上训练得到，当策略生成分布逐渐偏离奖励模型训练分布时，奖励评分的可靠性会下降，进一步导致奖励模型过优化\cite{gao2023scaling}。已有研究尝试通过限制策略偏移或更新反馈数据缓解这一问题\cite{moskovitz2023fool,online_rlhf_2025}，但这些工作主要从策略侧或数据侧入手，对奖励函数自身在跨轮更新中的变化幅度及其对整体对齐效果的影响讨论仍不充分。

\subsection{直接偏好优化及其变体}

为降低RLHF中在线强化学习训练的复杂度，Rafailov等人提出直接偏好优化（Direct Preference Optimization, DPO）\cite{rafailov2023direct}。DPO将偏好对数据直接转化为策略优化目标，省去了显式奖励模型训练和在线强化学习环节，因此实现流程更简单，训练开销也更低。后续相关方法也从目标函数形式、数据利用方式和参考模型约束等角度扩展了这一方向\cite{azar2024general,ethayarajh2024kto,meng2024simpo,gorbatovski2025trdpo}。

DPO系列方法的主要局限在于，偏好监督主要来自离线数据，奖励信号则以策略概率比的形式隐式表达。随着策略参数更新，新策略生成的样本未必仍落在原始偏好数据覆盖的范围内，因此长时间训练和分布外泛化仍可能受到离线数据分布限制\cite{tajwar2024preference}。从本文关注的问题看，DPO路线表明，弱化显式奖励建模能够简化偏好优化流程，但由于奖励信号隐含在策略目标中，其尺度和方向难以像显式奖励函数那样被单独检查、更新和约束。因此，DPO路线虽然简化了偏好优化流程，但并未直接回答奖励函数如何随策略分布变化而更新、以及更新幅度如何受控的问题。

\subsection{逆强化学习在语言模型对齐中的应用}

逆强化学习（Inverse Reinforcement Learning, IRL）由Ng和Russell\cite{ng2000algorithms}提出，其目标是从专家示范行为中反推潜在奖励函数。在语言模型对齐场景中，高质量示范响应可以被视为专家示范行为的一种序列化形式，也是偏好目标在行为层面的近似体现，因此IRL为从示范响应中恢复奖励信号提供了理论框架。最大熵逆强化学习\cite{ziebart2008maximum}进一步给出了常用的概率建模方式，使奖励函数与策略分布之间的关系能够被形式化刻画。与典型RLHF奖励建模依赖成对偏好比较不同，IRL更强调从示范行为本身学习奖励信号，这为减少额外偏好比较数据依赖提供了另一种思路。

近年来，IRL思想开始被用于语言模型后训练和偏好对齐。相关研究指出，IRL可以从示范行为中学习可用于策略优化的奖励函数，并为动态奖励学习提供建模基础\cite{irl_llm_survey_2025,wulfmeier2024irl}。也有工作从示范数据利用和奖励动态调整等角度推进了这一路线\cite{zeng2025d2r,cheng2026drirl}。这些研究说明，奖励函数不必始终作为固定打分器存在，也可以作为训练过程中持续更新的组件。

然而，一旦奖励函数参与多轮策略--奖励交替优化，新的稳定性问题也会出现。如果相邻轮次奖励函数变化过大，奖励尺度可能膨胀，奖励函数所刻画的偏好语义也可能发生漂移，进而影响策略优化方向。现有研究较少从双层优化角度刻画相邻轮次奖励变化与原始对齐目标改进之间的关系，也缺少可直接嵌入训练流程的奖励变化约束机制。本文后续提出的PIRO即围绕这一问题展开，重点研究逆强化学习框架下奖励函数跨轮变化的约束机制。

\section{论文主要工作内容}

本课题围绕基于逆强化学习的大模型偏好对齐方法展开，重点研究奖励—策略交替更新过程中如何约束奖励函数的跨轮变化，以提升动态奖励学习的稳定性。全文工作包括理论分析、算法设计和实验验证三个方面。

（1）\textbf{理论分析}：本文将大语言模型对齐建模为双层优化（bilevel optimization）问题。外层以奖励函数参数为优化变量，从高质量示范响应与当前策略样本的差异中学习奖励信号；内层以策略参数为优化变量，在奖励函数与策略偏移正则项的共同作用下更新语言模型策略。在该建模下，本文推导内层最优策略的 Boltzmann 形式，构造局部替代目标，并证明其在旧奖励参数处与原始目标具有一阶等价性。在此之上，本文给出改进下界定理，刻画相邻两轮奖励函数变化幅度对局部替代目标与原始对齐目标偏差的影响：当奖励函数更新过大时，替代目标对原始目标的近似误差会增大，即使替代目标在当前轮次得到优化，原始对齐目标仍可能无法稳定改进。这说明在IRL式奖励—策略交替训练中，必须对奖励函数更新幅度加以约束。

（2）\textbf{算法设计}：上述理论结论表明需要把相邻两轮奖励函数的最大变化量控制在一定范围内，但该范围取决于奖励取值范围、输出空间规模、策略正则强度和单轮目标改进量等因素。语言模型可能生成的输出序列数量极大，单轮真实改进量也无法在训练前预知，这使理论上的硬约束难以直接用于实际训练。本文引入相邻迭代轮次中奖励模型在同一批样本上输出差异的均方根差（Root Mean Square Difference, RMSD），作为奖励函数更新幅度的可观测代理量。基于该代理量，本文设计自适应约束系数调节机制：RMSD超过上阈值时增大约束系数以加强惩罚，低于下阈值时减小约束系数以避免约束过强。由此，理论中的奖励变化约束思想被转化为可嵌入训练循环的自适应反馈机制。

（3）\textbf{实验验证}：本文在摘要生成和通用指令跟随两个任务上验证PIRO。实验将PIRO与SFT、DPO、固定奖励模型下的RLHF(PPO/GRPO)、无约束IRL基线以及去除RMSD约束的消融版本进行比较。结果表明，PIRO在两类任务上均取得更稳定的训练过程和更好的综合评价表现；去除奖励更新约束后，模型更容易出现多轮训练退化。这说明RMSD约束能够缓解动态奖励学习中的不稳定问题，也支持了本文关于奖励函数更新幅度需要受控的理论判断。


\section{论文结构安排}

全文共分为五章，各章内容安排如下。

第一章为绪论。先介绍大语言模型对齐问题的研究背景，说明固定奖励模型范式下分布偏移与奖励模型过优化问题的来源；再综述RLHF、DPO和IRL在语言模型对齐中的研究进展，引出IRL框架下奖励函数动态更新的稳定性问题；最后概括本文提出的PIRO方法及其在理论分析、算法设计和实验验证三个方面的主要工作。

第二章为相关理论与技术基础。本章按策略优化、奖励建模、偏好优化和奖励函数学习的递进关系展开：先介绍强化学习基础以及PPO、GRPO等策略优化算法；接着说明RLHF如何用偏好数据训练奖励模型并驱动策略更新；再介绍DPO如何把显式奖励建模转化为离线偏好优化；最后讨论逆强化学习和最大熵逆强化学习如何从高质量示范响应中学习可更新的奖励函数。该章重点说明奖励建模、策略偏移约束、分布偏移和奖励模型过优化之间的关系，为第三章的双层优化建模和奖励函数更新约束推导做理论准备。

第三章为受约束逆强化学习对齐方法。本章把大语言模型对齐建模为奖励函数学习与策略优化耦合的双层优化问题，在最大熵IRL框架下推导内层策略的Boltzmann闭式形式；接着构造替代目标并证明其与原始目标的一阶等价性，由此给出改进下界定理，说明奖励函数更新幅度过大会增大局部替代目标与原始目标之间的偏差，进而影响对齐目标的稳定改进。在此之上，本章将理论中的奖励变化约束思想转化为可实现的 RMSD 代理正则项，设计基于RMSD的自适应约束系数调节机制，并给出PIRO的完整奖励—策略交替训练流程。

第四章为实验设计与结果分析。本章在摘要生成和通用指令跟随两个任务上评估PIRO，并与SFT、DPO、固定奖励模型RLHF、无约束IRL基线及去约束消融方法进行比较，重点验证RMSD约束对多轮IRL训练稳定性和通用对齐效果的影响。具体实验设置、评价指标和结果分析在该章展开。

第五章为总结与展望。本章总结全文主要结论，归纳本文在奖励函数更新约束理论、RMSD自适应调节机制和实证验证方面的贡献；并分析方法在实验规模、奖励信号来源和多轮在线数据收集上的局限，展望在更大模型、多维奖励建模和在线偏好学习融合方向上的扩展可能。

